\section{对象、原生组织与逻辑映射}

本例采用 HZO 二元 1T1C 存储阵列与局部数字乘加。极化状态经感测放大器（SA）判为 bit，再执行整数归约；破坏性读取之后恢复全部激活行。所选参考配置为 $K=N=128$ 的 INT8 矩阵，输入128 Byte、有效容量16384 Byte，输出128个23-bit整数容器。它由有实测依据的 HZO 存储机制构造，所述数字 CIM 组织属于工程参考配置。

\sid{FERAM-03} 的64 kbit阵列采用 TiN/HZO/TiN、11 nm HZO，在2.5 V、1\,$\mu$m$^2$电容条件下给出二元存储证据。本文保留该存储电压与负载锚点，28 nm、近0.9 V条件用于数字外围；不以膜厚代替工艺节点，也不按节点比例缩短极化时间。\cite{feram03}

每权重使用八个独立二元 FeCAP。令输入位置 $i$、输出位置 $j$、权重 bit $k$，则
\[
\text{分片}=i\bmod32,\quad
\text{局部行}=8\lfloor i/32\rfloor+\lfloor j/16\rfloor,\quad
\text{列}=8(j\bmod16)+k.
\]
32片各有32条128-bit局部行，共131072个数据cell；没有权重复制。该局部行长度提供足够短的独立感测、恢复及外部更新域。32片不是32个独立矩阵，而是同一矩阵的输入行分组。

\begin{table}[htbp]\centering\small
\begin{tabularx}{\textwidth}{@{}p{30mm}X@{}}\toprule
资源 & 参考组织\\\midrule
读与恢复 & 每片128 SA及恢复 BL 驱动，共4096路；32个局部 PL 域同时恢复32条完整行。每 SA 一条本地参考电容路径，共4096条，参考不计有效权重。\\
读码保持 & 单个4096-bit操作数寄存器，已有捕获包含于读窗口；固定组内跨八个输入bit只读复用，完成后覆盖。\\
数字计算 & 16输出通道，每轮32项、权重八位并行，输入逐bit；128个独立23-bit累加器跨行组保存部分和。无 ADC。\\
本地接口 & 输入寄存1024 bit，结果寄存2944 bit；捕获与最终提交各一个公共数字拍。\\
外部更新 & 单个更新域、128-bit编码接口、128 BL目标驱动、一条所选局部 PL；一笔更新16个完整INT8权重。\\\bottomrule
\end{tabularx}
\caption{原生局部资源账本。4096路内部恢复与128路外部目标写入分别由本地读码和外部数据控制，不能混同。}
\end{table}

每次选一组32输入项、16输出，共4096 bit，全部感测并恢复32条完整行。四个输入行组、八个输出组共有32次物理读；各组读码跨八个输入bit使用，计算256轮。操作数写使能在计算期间关闭，独立累加器接收结果；下一组捕获不与当前组重叠。因此无需第二个操作数 bank，也不能省去真实破坏读恢复。每向量读恢复覆盖131072个数据cell各一次。

\section{原始证据与完整服务参数}

\sid{FERAM-03} p.2 Figs.8--9显示 BL预充、与参考比较、SA放大引起状态翻转，以及 data writeback。p.1正文和Fig.11在2.5 V给出8 ns sense、14 ns write的无失败测试点。8 ns是感测窗口；14 ns原标为write latency，没有被分解为单独极性脉冲。Fig.11扫描1--50 ns；Figs.10、12采用100 ns操作条件。\cite{feram03}

令 $t_s$ 为感测及 SA 锁存，$t_p$ 为参考设计每极性保持，$t_e$ 为专用建立与关闭合计。$t_e$ 包含 BL预充/参考形成、WL/PL偏置与边沿、已有读码捕获、SA复位及返回可复用状态，排除 $t_s,t_p$。读侧一相破坏性恢复计为
\begin{equation}
t_{m,D}=t_s+t_p+t_e.
\label{08_feram_hfo2:eq:mediaread}
\end{equation}

\begin{table}[htbp]\centering\small
\begin{tabularx}{\textwidth}{@{}p{32mm}rrrX@{}}\toprule
时间（ns） & 短 & 参考 & 长 & 采用依据\\\midrule
$t_s$感测/锁存 & 8 & 20 & 50 & 8 ns实测锚点，20/50 ns为同一操作扫描内较长设计窗口。\\
$t_p$每极性保持 & 14 & 50 & 100 & 14 ns来自原write标签；为所选两相混合行更新，每相保留完整窗口。50/100 ns由扫描和表征条件支持量级。\\
$t_e$建立/关闭 & 20 & 40 & 80 & 前后各10/20/40 ns的工程预算；包含独立偏置边沿、捕获与复位。\\
完整读 & 42 & 110 & 230 & $t_s+t_p+t_e$，含一次整行恢复。\\
完整物理写 & 48 & 140 & 280 & $2t_p+t_e$，外部接口另计。\\\bottomrule
\end{tabularx}
\caption{同一组织、同一安装资源额定能力的成对服务条件。原值与工程采用窗口分开。}
\end{table}

外部混合0/1行采用两相 PL 组织：目标0列 BL保持0，目标1列保持 $V_W$；PL先保持 $V_W$ 再保持0。$V_{FE}=V_{PL}-V_{BL}$ 因而分别经历 $(+V_W,0)$ 和 $(0,-V_W)$，两种旧状态均能到达所需目标。两相各付 $t_p$，BL码跨两相保持，建立关闭只付一次；这是所选参考波形，不把 $2t_p$ 报为原文14 ns写测量。二元开环写不增加未报告的模拟verify、重试或擦除。

恢复窗口的量级由其他 HZO 阵列交叉核查。\sid{FERAM-02} 的5 nm HZO阵列在感测后写回、等待 $t_{WR}$ 再预充关断；20 ns约覆盖70\%极化响应，约65 ns附近信号趋于饱和。\sid{FERAM-04} 报告完整 $t_{RC}=185$ ns，Data 0恢复与预充重叠，$t_{RP}=80$ ns不能再作为独立纯边沿与一次完整恢复叠加。上述来源只支持完整阶段与量级，不替代本例不同膜厚/电压的成对参数。\cite{feram02,feram04}

\section{驱动负载、保持生命周期与条件}

\sid{FERAM-03} Fig.12为130 nm读裕量模型采用 $C_{BL}=250$ fF，40 nm估算为120 fF。本文保留250 fF和2.5 V作为声明的局部负载，未按缩小后的节点取更快值。每条 BL 在10/20/40 ns边沿内充放电的电容电流为
\[
I_{BL}=C_{BL}V_W/t_e^{\rm edge}=62.5/31.25/15.625\,\mu\mathrm{A}.
\]
固定配置每条 BL 预留80\,$\mu$A驱动资格，三情景不更换驱动。128条外部写 BL 的需求为8/4/2 mA；4096条并行恢复 BL 为256/128/64 mA。这些是相应边沿的电容负载需求，额定驱动及供电分布需支持最快条件，不能以典型4 mA作为总供给上限。

PL还承担铁电翻转。原文 $2P_r>40\,\mu\mathrm{C/cm^2}$ 是材料实测下界，配1\,$\mu$m$^2$面积给出每cell翻转电荷大于400 fC、每128-cell PL大于51.2 pC。以每相14/50/100 ns计，单PL极化电流下界为3.66/1.024/0.512 mA；32 PL并行对应大于117/32.8/16.4 mA。介电电荷和布线电容还需另计，因此这些下界不能作为充分的PL额定值。BL边沿与极化保持分阶段安排，不将不同相位的峰值直接相加。完整PL负载及并行恢复供电尚需按该构造组织确认，是主要电路资格条件。

SA判决后，4096-bit操作数先捕获，再完成恢复、关行与SA复位。该寄存器在八个计算拍期间保持，参考为40 ns，随后可被下一组覆盖；累加器维持到完整向量结束。其生命周期符合公共单tile政策；新增存储为零，捕获已在 $t_e$ 中覆盖，不重复收取32拍。整个过程按顺序服务，不使用恢复与计算重叠或全矩阵shadow。

\section{服务推导、结果及物理解释}

使用公共数字节拍 $T_D=2/5/10$ ns和显式逻辑形状，
\begin{equation}
\Delta_S=32(t_s+t_p+t_e)+(256+2)T_D.
\label{08_feram_hfo2:eq:stream}
\end{equation}
每笔完整更新一条128-bit局部行，即 $B_R=16$ Byte；编码首拍与命令同拍，准备/完成共 $2T_D$：
\begin{equation}
\Delta_R=2T_D+t_e+2t_p,\qquad
\rho=128/\Delta_S,\quad\tau=16/\Delta_R,\quad\RI^*=8\Delta_R/\Delta_S.
\label{08_feram_hfo2:eq:resident}
\end{equation}
纳秒代入得到十进制GB/s。32片各32行，共1024笔完整行事务覆盖矩阵一次，$T_R=1024\Delta_R$。小于整行的局部更新若要保留其余内容，须加入相应读保留服务，不能直接按字节缩短完整行时间。

\begingroup
\let\feramresulttable\table
\renewcommand{\table}[1][htbp]{\feramresulttable[H]}
% Generated by check_feram.py --emit.
\begin{table}[htbp]\centering\small
\begin{tabular}{@{}lrrrrrr@{}}\toprule
情景 & $\Delta_S$ ($\mu$s) & $\Delta_R$ (ns) & $\rho$ (GB/s) & $\tau$ (GB/s) & $\RI^*$ & 整矩阵写 ($\mu$s)\\\midrule
短 & 11.268 & 52 & 0.01136 & 0.30769 & 0.03692 & 53.248\\
参考 & 29.450 & 150 & 0.00435 & 0.10667 & 0.04075 & 153.600\\
长 & 61.460 & 300 & 0.00208 & 0.05333 & 0.03905 & 307.200\\
\bottomrule\end{tabular}
\caption{同一1T1C二元映射的成对工程情景。$B_S=128$ Byte、$B_R=16$ Byte；末列为1024次局部行事务的聚合。}
\end{table}
\begin{table}[htbp]\centering\small
\begin{tabular}{@{}lrrrrr@{}}\toprule
组织（参考时隙） & 读/计算轮 & $\Delta_S$ ($\mu$s) & $\rho$ (GB/s) & $\tau$ (GB/s) & $\RI^*$\\\midrule
D0 每输入bit重读 & 256/256 & 29.450 & 0.00435 & 0.10667 & 0.04075\\
D1 锁存后八次复用 & 32/256 & 4.970 & 0.02575 & 0.10667 & 0.24145\\
\bottomrule\end{tabular}
\caption{只改变数字锁存复用。D1另有4096-bit权重锁存和32拍捕获，不改resident更新宽度。}
\end{table}

\endgroup

典型 $\Delta_S=32(20+50+40)+258\times5=4810$ ns；$\Delta_R=10+40+100=150$ ns。于是 $\rho\approx27$ MB/s、$\tau\approx1.1\times10^2$ MB/s、$\RI^*\approx0.25$。完整装载 $T_R=153600$ ns，对应
\[
U^*=T_R/\Delta_S=128\RI^*\approx31.9\ \text{个向量}.
\]
这是与该矩阵装载和向量服务匹配的两表接口。N、输入共享或更新域改变时须重算，而不是直接沿用宏级阈值。

位置主要由两点决定。读码复用使每个cell每向量只承担一次破坏性恢复；典型恢复仍占1600/4810约33\%，建立关闭与数字计算各约27\%。写侧两极性保持占100/150约67\%，同时受到单外部更新域约束。两路共享 $t_p$ 和数字节拍，使成对RI*变化较小； $\rho$ 和 $\tau$ 本身跨三个正常窗口约有5--6倍差异。固定其他参数的 $t_p$ 对照保留主导窗口的独立影响。

同硬件逐输入bit重读的调度对照为 $256\times110+258\times5=29450$ ns，体现主动重复感测和恢复的成本；它不是1T1C必需调度。容量、资源与单次介质参数均不因该对照改变。

\FloatBarrier
\needspace{18\baselineskip}
\section{维护与适用范围}

主状态存于剩余极化，活动服务窗口不额外设置周期刷新；每次真实破坏读的恢复已含于完整读服务，原始与长期有效能力在此维护假设下相同。\sid{FERAM-03} 的85\,$^\circ$C保持实测为100 min，10年为投影；$10^{11}$耐久表述含频率外推。寿命属于可行性约束，不转换为瞬时 $\tau$。最坏高态cell每向量可承担一次读恢复极化循环，还需计外部更新。

\sid{FERAM-01} 的2T2C C2FeRAM通过隔离FeCAP端、耦合电容调制读晶体管，提供另一种非破坏读机制；其中一个电容是普通耦合电容，浮置节点泄漏仍可能要求restore。其离散电容实验与fF级模拟不构成本例1T1C阵列的周期或保持证据，故仅保留机制边界，不导入其数值。\cite{feram01}

本例的主要工程判断是两相保持窗口、独立BL/PL局部控制，以及4096路感测恢复的负载资格。结果以这些明确资源和状态判决条件成立为前提；输入、未取整结果、完整装载接口和原文定位保存在案例JSON，公共方法统一执行维度、payload与阶段计数。
